%%%%%%%%%%%%%%%%%%%%%%%%%%%%%%%%%%%%%%%%%
% Jacobs Landscape Poster
% LaTeX Template
% Version 1.0 (29/03/13)
%
% Created by:
% Computational Physics and Biophysics Group, Jacobs University
% https://teamwork.jacobs-university.de:8443/confluence/display/CoPandBiG/LaTeX+Poster
% 
% Further modified by:
% Nathaniel Johnston (nathaniel@njohnston.ca)
%
% This template has been downloaded from:
% http://www.LaTeXTemplates.com
%
% License:
% CC BY-NC-SA 3.0 (http://creativecommons.org/licenses/by-nc-sa/3.0/)
%
%%%%%%%%%%%%%%%%%%%%%%%%%%%%%%%%%%%%%%%%%

%----------------------------------------------------------------------------------------
%   PACKAGES AND OTHER DOCUMENT CONFIGURATIONS
%----------------------------------------------------------------------------------------

\documentclass[final]{beamer}

% Use the beamerposter package for laying out the poster
% Edit this scale to change the font size, and all rendered graphics
\usepackage[scale=1.8]{beamerposter}
\usepackage{import}

% Allow for larger equations
\usepackage{lmodern}
\usepackage{exscale} 

% Use the confposter theme supplied with this template
% Look in a sub folder: http://tex.stackexchange.com/a/284157
\usepackage{./packages/beamerthemeconfposter}

% Many more colors are available for use in beamerthemeconfposter.sty
\setbeamercolor{headline title}{fg=white,bg=RoyalBlue!95} % Colors of the block titles
\setbeamercolor{headline line}{fg=black,bg=Goldenrod} % Colors of the block titles
\setbeamercolor{block title}{fg=RoyalBlue,bg=white} % Colors of the block titles
\setbeamercolor{block body}{fg=black,bg=white} % Colors of the body of blocks
\setbeamercolor{block alerted title}{fg=black,bg=Goldenrod!70} % Colors of the highlighted block titles
\setbeamercolor{block alerted body}{fg=black,bg=Goldenrod!10} % Colors of the body of highlighted blocks

% set colors for itemize/enumerate
\setbeamercolor{item}{fg=Goldenrod}
\setbeamercolor{item projected}{fg=black,bg=Goldenrod}

%-----------------------------------------------------------
% Define the column widths and overall poster size
% To set effective sepwid, onecolwid and twocolwid values, first choose how many columns you want and how much separation you want between columns
% In this template, the separation width chosen is 0.024 of the paper width and a 4-column layout
% onecolwid should therefore be (1-(# of columns+1)*sepwid)/# of columns e.g. (1-(4+1)*0.024)/4 = 0.22
% Set twocolwid to be (2*onecolwid)+sepwid = 0.464
% Set threecolwid to be (3*onecolwid)+2*sepwid = 0.708

\newlength{\sepwid}
\newlength{\onecolwid}
\newlength{\twocolwid}
\newlength{\threecolwid}
\setlength{\paperwidth}{60in} % A0 width: 46.8in
\setlength{\paperheight}{40in} % A0 height: 33.1in
\setlength{\sepwid}{0.024\paperwidth} % Separation width (white space) between columns
\setlength{\onecolwid}{0.22\paperwidth} % Width of one column
\setlength{\twocolwid}{0.464\paperwidth} % Width of two columns
\setlength{\threecolwid}{0.708\paperwidth} % Width of three columns
\setlength{\topmargin}{-0.5in} % Reduce the top margin size
%-----------------------------------------------------------

\usepackage{graphicx}  % Required for including images
\usepackage{tikz} % Allow placing of logo in top right
\usepackage{booktabs} % Top and bottom rules for tables

%----------------------------------------------------------------------------------------
%   TITLE SECTION 
%----------------------------------------------------------------------------------------

\title{Unnecessarily Complicated Research Title} % Poster title

\author{John Smith, James Smith and Jane Smith} % Author(s)

\institute{University of Delaware - Department of Mechanical Engineering} % Institution(s)

%----------------------------------------------------------------------------------------

\begin{document}

\addtobeamertemplate{block end}{}{\vspace*{2ex}} % White space under blocks
\addtobeamertemplate{block alerted end}{}{\vspace*{2ex}} % White space under highlighted (alert) blocks

\setlength{\belowcaptionskip}{2ex} % White space under figures
\setlength\belowdisplayshortskip{2ex} % White space under equations

% The whole poster is enclosed in one beamer frame
\begin{frame}[t]

%----------------------------------------------------------------------------------------


% Add logo in top right
\tikz[remember picture,overlay]
  \node[anchor=north east,inner sep=0in,shift={(-0.7in,-1.1in)}] at (current page.north east)
    {\includegraphics[height=3in]{figures/udel_logo.png}};

%----------------------------------------------------------------------------------------


% Add some padding above the first row
\vspace{20mm}

% Begin the first row of blocks
\begin{columns}[t,totalwidth=\paperwidth]

    % Empty spacer column
    \begin{column}{\sepwid}\end{column}

    % The first block (block 1.1)
    \begin{column}{\onecolwid}\vspace{-.6in}
        %----------------------------------------------------------------------------------------
%   MATERIALS
%----------------------------------------------------------------------------------------

\begin{block}{Materials}

The following materials were required to complete the research:

\begin{itemize}
\item Curabitur pellentesque dignissim
\item Eu facilisis est tempus quis
\item Duis porta consequat lorem
\item Eu facilisis est tempus quis
\end{itemize}

The materials were prepared according to the steps outlined below:

\begin{enumerate}
\item Curabitur pellentesque dignissim
\item Eu facilisis est tempus quis
\item Duis porta consequat lorem
\item Curabitur pellentesque dignissim
\end{enumerate}

\end{block}
    \end{column} % End of column 2.1
    
    % The second block (block 1.2)
    \begin{column}{\onecolwid}\vspace{-.6in}
        \begin{block}{Methods}

Lorem ipsum dolor \textbf{sit amet}, consectetur adipiscing elit. Sed laoreet accumsan mattis. Integer sapien tellus, auctor ac blandit eget, sollicitudin vitae lorem. Praesent dictum tempor pulvinar. Suspendisse potenti. Sed tincidunt varius ipsum, et porta nulla suscipit et. Etiam congue bibendum felis, ac dictum augue cursus a. \textbf{Donec} magna eros, iaculis sit amet placerat quis, laoreet id est. In ut orci purus, interdum ornare nibh. Pellentesque pulvinar, nibh ac malesuada accumsan, urna nunc convallis tortor, ac vehicula nulla tellus eget nulla. Nullam lectus tortor, \textit{consequat tempor hendrerit} quis, vestibulum in diam. Maecenas sed diam augue.

\end{block}
    \end{column}
    
    % The third block (block 1.3)
    \begin{column}{\onecolwid}\vspace{-.6in}
        \begin{block}{Stuff}

Lorem ipsum dolor \textbf{sit amet}, consectetur adipiscing elit. Sed laoreet accumsan mattis. Integer sapien tellus, auctor ac blandit eget, sollicitudin vitae lorem. Praesent dictum tempor pulvinar. Suspendisse potenti. Sed tincidunt varius ipsum, et porta nulla suscipit et. Etiam congue bibendum felis, ac dictum augue cursus a. 

\end{block}


\begin{alertblock}{Important Result}

Lorem ipsum dolor \textbf{sit amet}, consectetur adipiscing elit. Sed commodo molestie porta. Sed ultrices scelerisque sapien ac commodo. Donec ut volutpat elit.

\end{alertblock}

    \end{column}
    
    % Empty spacer column
    \begin{column}{\sepwid}\end{column}

\end{columns}

% Force the content to be at the bottom of the page
\vfill

% Begin the second row of blocks
\begin{columns}[t,totalwidth=\paperwidth]

    % Empty spacer column
    \begin{column}{\sepwid}\end{column}

    % The first block (block 2.1)
    \begin{column}{\onecolwid}\vspace{-.6in}
        \begin{block}{Mathematical Section}

Nam quis odio enim, in molestie libero. Vivamus cursus mi at nulla elementum sollicitudin. Nam quis odio enim, in molestie libero. Vivamus cursus mi at nulla elementum sollicitudin.
  
\begin{equation}
E = mc^{2}
\label{eqn:Einstein}
\end{equation}

Nam quis odio enim, in molestie libero. Vivamus cursus mi at nulla elementum sollicitudin. Nam quis odio enim, in molestie libero. Vivamus cursus mi at nulla elementum sollicitudin.

\begin{equation}
\cos^3 \theta =\frac{1}{4}\cos\theta+\frac{3}{4}\cos 3\theta
\label{eq:refname}
\end{equation}

Nam quis odio enim, in molestie libero. Vivamus cursus mi at nulla elementum sollicitudin. Nam quis odio enim, in molestie libero. Vivamus cursus mi at nulla elementum sollicitudin.

\begin{equation}
\kappa =\frac{\xi}{E_{\mathrm{max}}} %\mathbb{ZNR}
\end{equation}

\end{block}
    \end{column} % End of column 2.1
    
    % The second block (block 2.2)
    \begin{column}{\onecolwid}\vspace{-.6in}
        \begin{block}{Results}

\begin{figure}
\includegraphics[width=0.8\linewidth]{figures/placeholder.jpg}
\caption{Figure caption}
\end{figure}

Nunc tempus venenatis facilisis. Curabitur suscipit consequat eros non porttitor. Sed a massa dolor, id ornare enim:


\begin{table}
\vspace{2ex}
\begin{tabular}{l l l}
\toprule
\textbf{Treatments} & \textbf{Response 1} & \textbf{Response 2}\\
\midrule
Treatment 1 & 0.0003262 & 0.562 \\
Treatment 2 & 0.0015681 & 0.910 \\
Treatment 3 & 0.0009271 & 0.296 \\
\bottomrule
\end{tabular}
\caption{Table caption}
\end{table}


\end{block}
    \end{column}
    
    
    % The third block (block 2.3)
    \begin{column}{\onecolwid}\vspace{-.6in}
        \begin{block}{Conclusion}

Nunc tempus venenatis facilisis. \textbf{Curabitur suscipit} consequat eros non porttitor. Sed a massa dolor, id ornare enim. Fusce quis massa dictum tortor \textbf{tincidunt mattis}. Donec quam est, lobortis quis pretium at, laoreet scelerisque lacus. Nam quis odio enim, in molestie libero. Vivamus cursus mi at \textit{nulla elementum sollicitudin}.

\end{block}
    \end{column}
    
    % The forth block (block 2.4)
    \begin{column}{\onecolwid}\vspace{-.6in}
        \begin{block}{Additional Information}

Maecenas ultricies feugiat velit non mattis. Fusce tempus arcu id ligula varius dictum. 
\begin{itemize}
\item Curabitur pellentesque dignissim
\item Eu facilisis est tempus quis
\item Duis porta consequat lorem
\end{itemize}

\end{block}
    \end{column}
    
    % Empty spacer column
    \begin{column}{\sepwid}\end{column}

\end{columns}


%----------------------------------------------------------------------------------------


\end{frame} % End of the enclosing frame


\end{document}
